    \begin{otherlanguage*}{english}
        This work investigates whether Big Five personality traits, induced
        through prompting in artificial intelligence agents based on large
        language models (LLMs), predict the investor profile of these
        agents, as measured by the IPA (Investor Profile Index) instrument.
        The question is broken down into three: whether LLMs consistently
        simulate the requested traits, which associations between
        personality and investor profile are robust, and how much they
        depend on the model and on generation parameters. In the main
        analysis, 300 agents were generated with the \texttt{llama3.1:8b}
        model, each with randomly assigned levels (Low/Medium/High) for the
        five traits, which answered the BFI-44 (personality self-report,
        used as a manipulation check) and the IPA. The analysis was
        repeated with other generation temperatures, with a new draw of
        personality profiles, and with other language models. Personality
        induction was faithful for all five traits ($\rho$ between 0.813
        and 0.929). In the main analysis, Openness to Experience,
        Extraversion, Conscientiousness, and Neuroticism significantly
        predicted the IPA score (OLS regression, R² = 0.459; p < 0.001),
        but only Openness to Experience and Extraversion remained
        significant across all temperatures, draws, and models tested; the
        Conscientiousness and Neuroticism effects proved unstable, and
        Agreeableness showed no consistent effect. The absolute investor
        profile assigned to the same personality varied strongly across
        models. The Openness to Experience finding converges with De
        Bortoli et al. (2019), Jiang, Peng and Yan (2024), and recent
        studies with LLMs. It is concluded that LLM-based AI agents, when
        properly induced and with verified manipulation, reproduce the
        direction of part of the known relationships between personality
        and risk tolerance, but that the level of risk they declare is
        largely a property of the model used.

        \vspace{\baselineskip}

        \noindent
        \textbf{Keywords}: Big Five. Investor profile. Artificial
        intelligence. Language models. Behavioral finance.
    \end{otherlanguage*}
